\subsection{Freezing the first support primes before a later third obstruction}
\label{sec:ghost-prefix-class3}

\begin{result}{C.1.20}{Proved}{The exact coefficient-five class-two prefix}
\label{lem:ghost-prefix-five}
Choose a first pair \(q_0,r_0\) by the original A.44
rules, and put
\[
 a_1=3q_0^{24},\qquad b_1=-2r_0^{24},\qquad
 n_j=(j+4)!.
                                                        \tag{45}
\]
Choose the old local roots \((\xi_p,\eta_p)\) in its
stage \(j=1\), so that
\(a_1\xi_p^{120}+b_1\eta_p^{120}=1\).
Let \(P,R\) be positive rational integers that are
units at \(q_0,r_0\), with
\[
 P\equiv\xi_p\pmod p,\qquad
 R\equiv\eta_p\pmod p
                \quad(p=q_0,r_0).                    \tag{46}
\]
There is a class-two homomorphism
\[
 \rho_*:G_{\mathbb Q}\longrightarrow
     H_5\rtimes\mathbb F_5^\times
                                                        \tag{47}
\]
with abelian Kummer pair \((a_1,b_1)\), uniquely
determined up to central geometric conjugacy by the
two inertia prescriptions in the proof. It is
unramified outside \(2,3,5,q_0,r_0\). Its image is
the full group of order five hundred, and its fixed
field \(L_*\) has maximal abelian subfield
\(\mathbb Q(\mu_5)\).

For every original A.44 continuation with the chosen
stage-one local roots and with
\[
 q_j\equiv r_j\equiv1\pmod{q_0r_0}
                       \quad(j\ge2),                 \tag{48}
\]
the coefficient-five reduction of every integral
class-two section, after aligning its abelian
cocycles, is centrally conjugate to \(\rho_*\).
In particular this field is fixed before the choice
of \(q_1\) and \(r_1\).
\end{result}
\begin{proof}
Fix integral five-adic Kummer cocycles
\(\alpha,\beta,p,r\) for \(a_1,b_1,P,R\),
respectively, and set
\[
 a^*=\alpha+120p,\qquad b^*=\beta+120r.
                                                        \tag{49}
\]
These are the actual Kummer cocycles of
\(a_1P^{120},b_1R^{120}\); their reductions are
literally \(\bar\alpha,\bar\beta\).
At \(q_0\) the second rational number is a principal
unit, by (46) and the old local equation. Hence
\(b^*=\delta u_0^*\) for an integral local
\(u_0^*\in\mathbb Z_5(1)\).
At \(r_0\) the first rational number is a principal
unit, so \(a^*=\delta v_0^*\).
The primitives are unique since
\(H^0(G_{\mathbb Q_p},\mathbb Z_5(1))=0\)
for these primes. The two prescriptions for a central
cochain \(h_*\), using the tame generators of
Proposition~\ref{prop:ghost-class3-residue}, are
\[
 h_*(T_{q_0})=24\bar u_0^*,\qquad
 h_*(T_{r_0})=-24\bar v_0^*.                           \tag{50}
\]
They are reductions of actual integral local class-two
sections: use the local cup primitives of
Proposition~\ref{prop:ghost-ordinary-cup}, with sign
chosen so that \(\delta h=-a\cup b\). Only the
vanishing of the principal-unit Kummer coordinate is
used, so (49) need not itself be an ordinary pair.

The finite cup \(\bar\alpha\cup\bar\beta\) is
zero globally. At each finite prime the original
A.44 invariant supplies a unit solution of
\(a_1X^{120}+b_1Y^{120}=1\); its two Kummer classes
modulo five are \(a_1,b_1\), so the local cup
vanishes. The finite global detection argument in
Proposition~\ref{prop:ghost-ordinary-cup} now applies.
Choose a normalized global primitive of its negative.
The residue surjection proved in the same proposition,
\[
 H^1(G_{\mathbb Q},\mathbb F_5(2))\longrightarrow
 \bigoplus_{p\ne5}H^0(\mathbb F_p,\mathbb F_5(1)),
                                                        \tag{51}
\]
adjusts that primitive to satisfy (50) and to be
unramified at all other primes outside
\(2,3,5,q_0,r_0\). The required residue vector has
finite support, since every finite-valued continuous
cochain factors through a finite Galois quotient.
At \(q_0,r_0\), compatibility of the prescribed
residues follows from the actual local primitives
just constructed. At two and three the residue group
in (51) is zero, so there is no further choice.
The resulting cochain gives (47) by the genuine
Heisenberg group law.

We record explicitly the uniqueness group:
\[
 H^1(G_{\mathbb Q,\{5,\infty\}},\mathbb F_5(2))=0.
                                                        \tag{52}
\]
Inflation injects its integral first-cohomology
version into
\(H^1(G_{\mathbb Q},\mathbb Z_5(2))=0\), by
Proposition~\ref{prop:ghost-h1-weight2}.
The integral second cohomology is
\(K_2(\mathbb Z)\{5\}=0\), by
\cite[VI.8.2 and Table VI.10.1.1]{GhostWeibelKbook},
as in the proof of
Theorem~\ref{thm:ghost-first-integral-class3}.
The coefficient sequence for multiplication by five
therefore proves (52). Any two choices of \(h_*\)
have difference a weight-two cocycle. It has zero
inertia at \(q_0,r_0\), by (50), and outside the
displayed ramification set by construction. At two
and three its tame relation is
\((p^2-p)w(T_p)=0\), whose scalar is invertible
modulo five. Thus it factors through the group in
(52), and is a coboundary. This is precisely central
geometric conjugacy in (47).

We next prove that every stated continuation gives
this same finite class, retaining the integral
primitives rather than just the finite abelian
classes. Choose integral Kummer cocycles for its
fresh primes and form the uniformly convergent sums
\[
 \begin{split}
 a&=\alpha+120\kappa(q_1)+
                   \sum_{j\ge2}n_j\kappa(q_j),\\
 b&=\beta+120\kappa(r_1)+
                   \sum_{j\ge2}n_j\kappa(r_j).
 \end{split}                                         \tag{53}
\]
Uniform convergence follows from
\(v_5(n_j)\longrightarrow\infty\), since all
the cocycle values are integral. Their cohomology
classes are exactly the given completed pair.
After an abelian geometric conjugation, an arbitrary
integral class-two section has these literal two
cocycles.

At either of \(q_0,r_0\), the quotient
\(q_1/P\) and every \(q_j\), \(j\ge2\), are
principal units; the corresponding assertion holds
for \(r_1/R\) and \(r_j\). The former congruences
are the stage-one CRT conditions (46), and the
latter are (48). Therefore
\[
 a-a^*=5\delta e_p,\qquad b-b^*=5\delta f_p
            \quad\text{on }G_{\mathbb Q_p},
                   \quad p=q_0,r_0,                  \tag{54}
\]
with \(e_p,f_p\in\mathbb Z_5(1)\).
For completeness, divide (53) minus (49) by five.
Each summand is an integral local coboundary because
its rational argument is a principal unit. Choose
its integral zero-cochain primitive. The coefficients
\(n_j/5\) tend to zero five-adically, so their
primitive series converges and proves (54) as an
identity of continuous cochains.

At \(q_0\), (54) gives
\(b=\delta(u_0^*+5f_{q_0})\); the actual tame
class-two relation therefore gives
\(h(T_{q_0})=24(u_0^*+5f_{q_0})\).
At \(r_0\) it gives
\(h(T_{r_0})=-24(v_0^*+5e_{r_0})\).
These are exactly (50) after reduction.
At every later support prime the valuation is
\(n_j\), divisible by five, and the same integral
relation makes its central inertia zero modulo five.
At a zero-valuation prime integral torsion-freeness
gives zero central inertia before reduction.
Thus \(\bar h-h_*\) is unramified outside five,
including at two and three by the scalar calculation
above. Equation (52) proves the required central
conjugacy. In particular no division by \(p-1\)
has been performed after reduction: the entire
freezing assertion follows from (54).

Finally, the two Kummer classes \(a_1,b_1\) are
independent over \(\mathbb Q(\mu_5)\), by their
valuations one at three and two, respectively.
The Heisenberg commutator therefore forces full
image \(H_5\) over that field. The cyclotomic
character adds \(\mathbb F_5^\times\), so the
image in (47) has order five hundred. Its commutator
subgroup is \(H_5\), exactly as in
Theorem~\ref{thm:ghost-first-integral-class3}.
This proves the asserted maximal abelian subfield
and completes the lemma.
\end{proof}

\begin{result}{C.2.7}{Proved}{First class-three failure inside the compact-passing original family}
\label{thm:ghost-passing-first-class3}
There is a recursively specified original A.44 output
\[
 A=3\prod_{j\ge0}q_j^{(j+4)!},\qquad
 B=-2\prod_{j\ge0}r_j^{(j+4)!}
                                                        \tag{55}
\]
for which every \(q_j\) has three a noncube and
every \(r_j\) has two a noncube. It therefore
passes every compact elliptic support-order test
(16) and admits the full profinite affine-cover
abelian section of
Proposition~\ref{prop:ghost-all-passing-genjac},
with proper class zero. Nevertheless its first
empty unrestricted integral nilpotent section fibre
is class three. All twenty-four integral class-two
section classes have a nonzero local third
obstruction at \(q_1\), at coefficient prime five,
already nonzero after reduction modulo twenty-five.
\end{result}
\begin{proof}
Choose \(q_0,r_0\) by the first stage of
Proposition~\ref{prop:ghost-elliptic-passing},
including its two noncube conditions. Retain (45).
At stage one, whose exponents are \(n=120,m=720\),
choose every old local root required by A.44.
At five impose \(\xi_5=1\). This is allowed
because \(b_1\) is a five-adic unit and
\[
 v_5(a_1+b_1-1)\ge H_1(5)=6,
                    \qquad v_5(120)=1.
\]
The original Hensel calculation gives a root
\(\eta_5\) with \(v_5(\eta_5-1)\ge5\),
which satisfies the required additional congruence
modulo \(5^{v_5(720)}\). All other roots retain
the original requirements. Form the original CRT
modulus \(M\) and class \(Q_0\) of A.44.15.
In particular
\[
 Q_0\equiv1\pmod{720},\qquad
                  Q_0\equiv1\pmod{25}.              \tag{56}
\]
The latter follows because the exponent of five
in \(M\) is \(H_2(5)=7\) and \(\xi_5=1\).
Choose integer placeholders \(P,R\) satisfying
(46), and construct the field \(L_*\) of
Lemma~\ref{lem:ghost-prefix-five} before choosing
the new primes.

Here is the exact finite field intersection for
the next prime. Put
\[
 \begin{split}
 F&=\mathbb Q(\mu_M),\qquad
 R_b=F(b_1^{1/120}),\\
 R&=F(b_1^{1/120},a_1^{1/5}),\qquad
 C=F(3^{1/3}).
 \end{split}                                         \tag{57}
\]
In (57), the field denoted \(R\) replaces the
integer placeholder of (46); that placeholder
will no longer be used in the field argument.
Let \(\sigma_F\) be the automorphism of \(F\)
specified by the CRT class \(Q_0\).
The original A.44 field-intersection argument
supplies an automorphism of \(R_b\) restricting
to \(\sigma_F\) and fixing the radical
\(b_1^{1/120}\).
The classes of \(a_1,b_1\) are independent in
\(F^\times/(F^\times)^5\). Indeed their rational
classes are independent by valuations at three
  and two; rational fifth-power classes inject into
a cyclotomic field. Indeed a nontrivial rational
class has a prime valuation not divisible by five,
forcing its radical degree to be five. A Galois
radical extension of that degree would contain
\(\mu_5\), of degree four, which is impossible.
Such a radical field therefore cannot be a subfield
of an abelian extension of \(\mathbb Q\).
The cyclic extension \(R_b/F\) has a unique
degree-five subextension, namely \(F(b_1^{1/5})\).
It follows that
\(R_b\cap F(a_1^{1/5})=F\). The automorphism
therefore extends to \(R\), fixing both displayed
radicals. Denote it by \(\sigma_R\).

The field \(R/F\) is abelian. By the maximal
abelian-subfield assertion of the lemma,
\(L_*\cap F=\mathbb Q(\mu_5)\), so
\(\operatorname{Gal}(L_*F/F)=H_5\).
Its maximal abelian subextension is
\(F(a_1^{1/5},b_1^{1/5})\), which is contained
in \(R\). Consequently the precise intersection is
\[
 L_*F\cap R=F(a_1^{1/5},b_1^{1/5}),\qquad
                        [L_*R:R]=5.                 \tag{58}
\]
Choose any nonzero central element \(z\in H_5\).
It fixes both fifth roots. Since \(\sigma_F\)
fixes \(\mu_5\), the pair of actions \(z\) on
\(L_*\) and \(\sigma_F\) on \(F\) gives an
automorphism of \(L_*F\). On the intersection
in (58) it agrees with \(\sigma_R\), so these
actions extend to \(L_*R\).

The additional noncube condition is compatible with
this entire compositum. The classes of three and
\(b_1\) are independent modulo cubes over \(F\),
by their rational valuations and the cyclotomic
injectivity (22). The only degree-three subfield
of \(R_b/F\) is \(F(b_1^{1/3})\); adjoining
\(a_1^{1/5}\) has degree five. Thus \(C\cap R=F\).
Equation (58), whose extra degree is five, now
implies \(C\cap L_*R=F\): otherwise
\([CR:R]=3\) would divide \([L_*R:R]=5\).
Choose on \(C\) the automorphism restricting to
\(\sigma_F\) and sending \(3^{1/3}\) to
\(\zeta_3 3^{1/3}\). The two actions extend to
the compositum \(L_*RC\).

Chebotarev's theorem
\cite[Theorem 8.31]{GhostMilneANT} supplies a prime
\(q_1\), beyond all the original bounds, with
this conjugacy class. It satisfies the original
CRT and radical conditions, and also
\[
 \begin{gathered}
 q_1\equiv1\pmod{25},\qquad
 3\notin(\mathbb F_{q_1}^\times)^3,\\
 \rho_*(\operatorname{Frob}_{q_1})=(0,0,c;1),
                         \qquad c\in\mathbb F_5^\times.
 \end{gathered}                                      \tag{59}
\]
These properties are invariant under conjugation:
the CRT action is in an abelian field, both radicals
are fixed together with their roots of unity, and
conjugation multiplies the two nonzero Kummer or
central coordinates by units. Choose the least
eligible prime. Then choose \(r_1\) by A.44.17
with the additional noncube condition of
Proposition~\ref{prop:ghost-elliptic-passing}.
Its old-prime CRT residues are exactly the chosen
\(\eta_p\), as required for the prefix lemma.

At all stages \(j\ge2\), impose (48) through
the choice of the old local roots at \(q_0,r_0\).
After stage one the coefficient sum is one modulo
both primes. When such a prime does not divide
\(n_j\), the simple-root Hensel argument gives
both old variables one modulo that prime, by
fixing one variable to one and solving for the
unit-coefficient variable near one. When it divides
\(n_j\), the original invariant
\(H_j(p)>2v_p(n_j)\) gives the same near-one
choice. When it divides \(n_{j+1}\), the original
stronger choice A.44.14 is retained as well.
Thus (48) is available indefinitely. The noncube
field intersections of
Proposition~\ref{prop:ghost-elliptic-passing}
apply to each resulting legal CRT class. Choose
all later primes with those noncube conditions.
No other A.44 condition is altered, so its proved
induction gives the ordinary infinite-support pair
(55), with every compact order test satisfied.

Let \(s\) be any integral class-two section above
this pair. The prefix lemma identifies its reduction
modulo five with \(\rho_*\), up to a central
coboundary after abelian alignment. At
\(q=q_1\) the coboundary is zero on Frobenius
modulo five, since \(q\equiv1\pmod5\).
Consequently its actual integral Frobenius
coordinates satisfy
\[
 \alpha=a(F_q)\equiv0\pmod5,\qquad
 \gamma=h(F_q)\equiv c\ne0\pmod5.                    \tag{60}
\]
The valuation vector is \((120,0)\).
Ordinarity gives \(b=\delta u\) locally, and
the exact tame relation gives
\(\beta=(q-1)u\), \(h(T_q)=120u\).
The first third-obstruction residue of
Proposition~\ref{prop:ghost-class3-residue} is therefore
\[
 Z_{D,q}=120(\gamma-q\alpha u)
          +(q^2-1)u\frac{120\cdot119}{2}.             \tag{61}
\]
By (59)--(60), its five-adic valuation is exactly
one, whereas \(v_5(q^2-1)\ge2\). Its local
integral cohomology class is nonzero, and its
reduction modulo twenty-five is nonzero. This
argument applies to all integral class-two sections;
it uses no extra finite-coefficient section classes.
Corollaries~\ref{cor:ghost-class2-lift} and
\ref{cor:ghost-class2-fibre24} give their existence
and exact number twenty-four, proving the first
failure assertion. The affine abelian cover section
is the one constructed for the entire compact-passing
domain in Proposition~\ref{prop:ghost-all-passing-genjac}.
\end{proof}

The finite field in (47) can be specified without
an unspecified fundamental-group computation. Enumerate
monic integral polynomials and certified Galois
realizations of \(H_5\rtimes\mathbb F_5^\times\),
with the cyclotomic quotient and the two specified
Kummer subextensions. Test ramification outside
\(2,3,5,q_0,r_0\) and the two tame inertia values
(50), using the chosen embeddings and generators.
The lemma proves that such a realization exists and
that every passing realization has the same kernel
up to the central marking, so this search terminates
and specifies \(L_*\). The two primitive residues
in (50) are finite arithmetic data: at \(q_0\),
compute \(b^*(F_{q_0})\) modulo
\(5^{1+v_5(q_0-1)}\) in its finite Kummer field
and divide the integral value by \(q_0-1\);
this determines \(u_0^*\bmod5\). Use \(a^*\)
at \(r_0\) in the same way. Integrality is
already proved by the principal-unit argument.
Its Frobenius test (59) is
then a finite number-field calculation. Combined
with least eligible prime and residue choices,
this makes (55) a recursive specification; no
numerical value for the enormous recursive primes
is asserted. The residue calculation (61) is the
same (m=120) calculation already checked in
\path{code/check_class3_first_example.py}.

\begin{result}{C.3.11}{Proved}{An affine abelian section with no metabelian cover lift}
\label{cor:ghost-metabelian-cover-exclusion}
For the pair (55), the fibre of sections of
\(\pi_1(Y)/\overline{[K,K]}\) above the original
completed coordinates is nonempty. The corresponding
fibre for
\[
 \pi_1(Y)/\overline{[[K,K],[K,K]]}
                                                        \tag{62}
\]
is empty. Thus every section in the former fibre
has no lift through the displayed metabelian
cover quotient.
\end{result}
\begin{proof}
Write \(K'=\overline{[K,K]}\) and
\(K''=\overline{[K',K']}\). Since
\(K\subset\Pi=\pi_1(U_{\overline{\mathbb Q}})\),
we have \(K'\subset\Gamma_2\Pi\) and hence
\[
 K''\subset\overline{[\Gamma_2\Pi,\Gamma_2\Pi]}
                     \subset\Gamma_4\Pi.
\]
The actual cover inclusion therefore induces a
continuous arithmetic homomorphism
\[
 \pi_1(Y)/K''\longrightarrow\pi_1(U)/\Gamma_4\Pi,
                                                        \tag{63}
\]
commuting with projection to \(G_{\mathbb Q}\)
and to the two original Kummer coordinates.
A section in the asserted source fibre would push
forward to an integral class-three section over
(55), contrary to the theorem. Nonemptiness of
the abelian fibre was proved there. The map
\(\pi_1(Y)/K''\to\pi_1(Y)/K'\) is the actual
quotient map, so the empty source fibre proves
the last assertion for each of its target sections.
\end{proof}
